\begin{theorem}[G-27: Nonrectangular boundary comparison through $2m-1$ rows]
\label{G-27}
Let $n>m\ge17$, $m<r\le2m-1$, and let
$\nu\vdash ns$, $s\ge1$, have exactly $r$ positive parts
and be nonrectangular. Put $a=n-m$. If $\nu_1\le as$,
the determinant multiplicity is strictly larger than the
padded permanent multiplicity at every sufficiently large
integer stretch. On the equality face the explicit threshold
below applies. Below the face the padded multiplicity is zero
at every positive stretch, and the determinant occurs at every
stretch beyond its explicit invariant-generator bound squared.

Together with G-26 and G-25, this classifies the eventual sign
in every fixed nonrectangular full-row direction with
$4\le r\le2m-1$. A positive permanent-minus-determinant gap
occurs eventually exactly when
\[
 \nu_1>(n-m)s,\qquad
 \Delta=rm^2-(r-2)n^2-2m+r-1>0.
\]
In every other case the determinant eventually dominates.
The strict-direction assertion already extends through $2m$
rows by G-25; the additional boundary assertion here stops
at $2m-1$ rows and excludes rectangular directions beyond
the range of G-26.
\end{theorem}

\paragraph{Explicit equality-face constants.}
Assume $\nu_1=as$. Put
\[
 q=r-1,\quad \mu=(\nu_2,\ldots,\nu_r),\quad
 \eta=\nu-\nu_r\mathbf1_r,\quad
 \theta=\mu-\nu_r\mathbf1_q,
 \quad e=|\eta|,\quad e'=|\theta|.
\]
As before, $r_\eta$ and $r_\theta$ count strict comparisons
between ordered parts. Define
\[
 \begin{aligned}
 u&=(r-2)n^2-r^2+3+r_\eta,\\
 v&=(r-1)m^2-2m-r^2+2r+2+r_\theta,\\
 b_D&=\binom{\binom{n+r-1}{r-1}+s-1}{s}
                    \dim S_\eta\mathbb C^r,\\
 b_B&=\binom{\binom{m+q-1}{q-1}+s-1}{s}
                    \dim S_\theta\mathbb C^q,\\
 b&=\max\{b_D,b_B\},\qquad
 L=ns+(r-1)e,\qquad p=ns+e,\\
 \beta&=\left\lceil\frac{3r^2b}{8}\right\rceil
                      L^{2(r^2-1)},\\
 T&=\binom{b+\beta}{b}-1,\qquad E=rT\beta!.
 \end{aligned}
\]
Then $u>v\ge2$, and strict determinant dominance holds for
every
\begin{equation}\label{G27:threshold}
                    j\ge E\bigl(4v(pE)^{v-1}+1\bigr)
\end{equation}
in coordinate degree $sj$. These are the same form of bounds
as G-26, but the lower invariant algebra is now allowed to
have a nontrivial degree period.

\begin{proof}
We compare the two enhanced invariant algebras and retain the
degree period only on the side used for an upper bound.
\paragraph{1. Separate stable algebras, with period one needed only above.}
On the equality face, the all-degree Pieri bound G-26 gives
\[
 m_{j\nu}(\ell^a\operatorname{per}_m)
       \le m_{j\mu}(\mathcal P_{m,q}).
\]
Let $A$ be the enhanced invariant algebra for
$\mathcal D_{n,r}$, outer degree $s$, and auxiliary weight
$\eta$, and let $B$ be that for $\mathcal P_{m,q}$ with
auxiliary weight $\theta$. Exact extraction yields
\[
 \dim A_j=m_{j\nu}(\det_n),\qquad
 \dim B_j=m_{j\mu}(\mathcal P_{m,q}).
\]
The two applications of R-32 are valid because
\[
 4\le r\le2m-1<2n-2,\qquad
 4\le q=r-1\le2m-2.
\]
Thus both enhanced cones have stable points with finite
projective stabilizers. Since $\nu$ is nonrectangular, the
first cone additionally has a stable point with no scalar
period. No such assertion is needed for $B$: the tail $\mu$
may be rectangular.

The exact image dimensions from R-28 and G-21 and the stable
quotient formula give $\dim A=u$ and $\dim B=v$.
The degree semigroup of $A$ has gcd one. The corresponding
gcd for $B$ is left unrestricted. This distinction is the
reason that no assertion about generic freeness of a
rectangular two-product tensor is required.

\paragraph{2. The dimension inequality persists beyond $m$ rows.}
Deleting the first part removes
\[
 r_\eta-r_\theta=\#\{j>1:\nu_1>\nu_j\}\ge0
\]
strict flag comparisons. Since $n\ge m+1$,
\[
 \begin{aligned}
 u-v
  &=(r-2)n^2-(r-1)m^2+2m-2r+1+r_\eta-r_\theta\\
  &\ge(2m-1)r-m^2-2m-1\\
  &\ge m^2-m-2>0.
 \end{aligned}
\]
For the last inequality use $r\ge m+1$. Moreover, writing
$q=r-1$, with $m\le q\le2m-2$, gives
\[
 v=q(m^2-q)-2m+3+r_\theta
        \ge m^3-2m^2+3>2.
\]
In particular $u\ge v+1$.

\paragraph{3. An upper Hilbert bound does not require degree period one.}
The degree and dimension estimates in G-23/G-26 give
homogeneous generator degrees at most $\beta$ in both rings,
at most $T$ generators, and standard $E$th Veroneses.
They also give $E\ge\beta^2$. The determinant no-scaling
point implies $A_j\ne0$ for every $j\ge\beta^2$.

For the lower algebra $B$, let $g$ be the gcd of the positive
degrees of a set of nonzero homogeneous generators. This
set is nonempty because $v\ge2$. Its degrees are at most
$\beta$. Divide them by $g$ and apply the elementary
residue-semigroup bound of G-18. It follows that every
multiple of $g$ at least $\beta^2$ belongs to their semigroup.
Since $B$ is a domain, a product of nonzero generators is
nonzero; hence $B_h\ne0$ at every such degree $h$.
Conversely $B_j\ne0$ implies $g\mid j$, by homogeneity of
the generators. This argument makes no saturation claim in
degrees not divisible by $g$.

Write $j=Ek+t$, $0\le t<E$. Since $g\le\beta$ and
$\beta!\mid E$, one has $g\mid E$. If $B_j\ne0$, then
$g\mid t$, and $2E-t\ge E\ge\beta^2$ is a multiple of
$g$. Multiplication by a nonzero element of $B_{2E-t}$
therefore embeds $B_j$ in $B_{E(k+2)}$. If $B_j=0$, the
same upper inequality holds without a multiplication map.
For $A$, multiplication by any nonzero element of
$A_{E+t}$ embeds $A_{E(k-1)}$ in $A_j$. Thus, for $k\ge1$,
\begin{equation}\label{G27:one_sided_sandwich}
 \dim A^{(E)}_{k-1}\le\dim A_j,\qquad
 \dim B_j\le\dim B^{(E)}_{k+2}.
\end{equation}
Only the lower bound in the first ring requires period one.

The lower permanent enhanced cone has a dense homogeneous
parameter map of degree $ms+e'\le p$. The image-degree
estimate used in G-23 therefore gives
\[
             \deg\operatorname{Proj}B^{(E)}\le(pE)^{v-1}.
\]
Linear Noether normalization for the standard domain
$A^{(E)}$ and the same Chardin bound for $B^{(E)}$ now give
\[
 \dim A_j\ge\binom{k+v-1}{v},\qquad
 \dim B_j\le(pE)^{v-1}\binom{k+v+1}{v-1}.
\]
As in G-26, their binomial ratio is at least $k/(4v)$ for
$k\ge v+1$, and \eqref{G27:threshold} makes the first
quantity strictly larger than the second. The Pieri upper
bound proves the desired strict comparison at every such
integer stretch, including all possible lower-ring residues.

\paragraph{4. The remaining two half-spaces.}
If $\nu_1<as$, G-26 gives zero padded multiplicity at every
positive stretch. R-32 still applies to the degree-$n$
determinant cone, because $r<2n-2$ and $\nu$ is nonrectangular.
It gives determinant occurrence at every
$j\ge\beta_D^2$, where
\[
 \beta_D=\left\lceil\frac{3r^2b_D}{8}\right\rceil
             \bigl(ns+(r-1)|\eta|\bigr)^{2(r^2-1)}.
\]
Here $b_D$ and $\eta$ have their displayed determinant
definitions; no lower tail degree is used in this case.
If $\nu_1>as$, G-25 supplies the comparison, with its
nonzero integer exponent $\Delta$. Combining these cases
with G-26 proves the stated eventual classification.
\end{proof}

\paragraph{Scope.}
The new semigroup argument allows a period on the algebra
used only for an upper bound; it does not remove that period.
Rectangular directions with more than $m$ rows, equality
faces at $2m$ rows, and strict directions beyond $2m$ rows
retain their separate stable-point gates. Earlier individual
degrees and varying-direction complexity questions are also
unsettled. The construction is independent of the stipulated
analytic inputs and makes no efficient-size claim for the
displayed thresholds.
